\chapter{Feedback: State-Space and Frequency-Domain Views}
\label{ch:feedback}
\Cref{ch:linsys,ch:splane} considers open-loop systems: an input drives an
output, but the output does not affect the input. Feedback measures the output,
compares it with the reference, and uses the difference to drive the system. That
operation changes the closed-loop poles rather than merely combining transfer
functions.

Like the rest of Part~II, this chapter stays generic. In the time and state-space
view, feedback changes differential-equation coefficients. In the frequency view,
it sets loop gain and stability margins. \Cref{ch:active} applies both views to
op-amps, oscillators, and regulators.

\section{Closing the Loop}
Let the system we want to control have transfer function \(A(s)\): in control
language, the \emph{plant}; in amplifier language, the \emph{forward path} or
\emph{open-loop gain}. Measure the output \(y\), scale it by a constant \(\beta\),
and subtract that from the reference input \(u\). The signal actually reaching the
plant is the \textbf{error} \(e=u-\beta y\), so
\[
Y=A(s)\,E=A(s)\bigl(U-\beta Y\bigr)
\quad\Longrightarrow\quad
\bigl(1+A(s)\beta\bigr)Y=A(s)\,U,
\]
which rearranges to the two definitions this chapter turns on:
\[
\boxed{\;
A_{\mathrm{cl}}(s)=\frac{Y(s)}{U(s)}=\frac{A(s)}{1+A(s)\beta},
\qquad
L(s)=A(s)\,\beta .\;}
\]
\(A_{\mathrm{cl}}\) is the \textbf{closed-loop} transfer function and \(L\) is the
\textbf{loop gain}, the gain accumulated once around the loop with it cut open.
This is the same result as the block-diagram rule \(G_{\mathrm{cl}}=G/(1+GH)\) in
\cref{ch:linsys}, renamed for amplifier work: \(G\to A\) for the forward path and
\(H\to\beta\) for the feedback path. The amplifier symbols are used from here on,
because they are the ones the radio literature and \cref{ch:active} use.

Because \(\beta\) is a constant, this is \textbf{proportional feedback}: the
feedback path scales the measured output without adding dynamics. Every result in
this chapter concerns that P-only model.

\begin{controlsbox}
\textbf{Where the poles come from.} The closed-loop poles are the roots of
\(1+L(s)=0\), i.e.\ \(1+A(s)\beta=0\). This is the single most useful equation in
classical control: it says that choosing \(\beta\) chooses where the poles go. The
root locus of \cref{ch:splane} is nothing more than the picture of this equation's
roots as \(\beta\) sweeps.
\end{controlsbox}

\section{The Large-Loop-Gain Limit}
When the loop gain magnitude is large at the frequency of interest,
\(|L|=|A\beta|\gg1\), the \(1\) in the denominator is negligible:
\[
A_{\mathrm{cl}}=\frac{A}{1+A\beta}\;\approx\;\frac{A}{A\beta}=\frac{1}{\beta}.
\]
The closed-loop behavior is then set by the feedback path \(\beta\) \emph{alone},
almost independently of \(A\). This is why feedback is the foundation of precision
engineering: \(A\) may be large, poorly known, temperature dependent, and
nonlinear, while \(\beta\) can be a precise, stable ratio. You spend surplus gain
to buy accuracy.

Quantitatively, differentiating \(A_{\mathrm{cl}}\) with respect to \(A\) gives the
\emph{sensitivity} result: a fractional change \(\dd A/A\) in the forward path
produces only
\[
\frac{\dd A_{\mathrm{cl}}}{A_{\mathrm{cl}}}
=\frac{1}{1+L}\,\frac{\dd A}{A},
\]
so the loop divides the plant's imperfections by \(1+L\).

\section{Proportional Feedback on a First-Order System}
Take the first-order plant of \cref{ch:linsys}. Its \textbf{DC gain} is its
zero-frequency gain, \(A(0)=A_0\), and its open-loop pole is at
\(-\omega_a\):
\[
A(s)=\frac{A_0}{1+s/\omega_a} .
\]
Close the loop with constant \(\beta\) and simplify:
\[
A_{\mathrm{cl}}(s)
=\frac{\dfrac{A_0}{1+s/\omega_a}}{1+\dfrac{A_0\beta}{1+s/\omega_a}}
=\frac{A_0}{1+s/\omega_a+A_0\beta}
=\underbrace{\frac{A_0}{1+A_0\beta}}_{\text{DC gain}}\cdot
\frac{1}{1+\dfrac{s}{\omega_a(1+A_0\beta)}} .
\]
The formula shows two linked effects. First, the closed-loop DC gain is
\(A_{\mathrm{cl}}(0)=A_0/(1+A_0\beta)\). Second, the pole moves from
\(-\omega_a\) to
\[
s_{\mathrm{cl}}=-\omega_a\,(1+A_0\beta),
\]
i.e.\ \emph{further left}, by the same factor. Left means faster: the closed-loop
time constant shrinks to \(\tau_{\mathrm{cl}}=\tau/(1+A_0\beta)\) and the bandwidth
widens by \((1+A_0\beta)\). Feedback trades gain for speed. The
\emph{gain--bandwidth trade} has product
\[
\bigl(\text{DC gain}\bigr)\times\bigl(\text{bandwidth}\bigr)
=\frac{A_0}{1+A_0\beta}\cdot\omega_a(1+A_0\beta)=A_0\,\omega_a
\]
is invariant. In \cref{fig:fb-first-order}, the left panel traces the pole moving
left as \(\beta\) increases; the right panel shows the corresponding faster,
lower-amplitude step responses.

\begin{figure}[htbp]
\centering
\includegraphics[width=0.78\textwidth]{feedback_first_order}
\caption{Proportional feedback on a first-order system with \(A_0=10\),
\(\omega_a=1\), and \(\beta=0,0.1,0.3,0.9\). \emph{Left:} the single
closed-loop pole \(-\omega_a(1+A_0\beta)\) slides left along the real axis as
\(\beta\) grows---a one-branch root locus that can never leave the real axis.
\emph{Right:} the step response gets faster and smaller because both its DC gain
and time constant decrease by the factor \(1+A_0\beta\).}
\label{fig:fb-first-order}
\end{figure}

\begin{physicalbox}
\textbf{Does proportional feedback make a first-order system second order?}
\emph{No.} The closed-loop denominator \(1+s/\omega_a+A_0\beta\) is still
first-degree in \(s\), so there is still exactly one pole and one state. For this
ideal first-order plant with positive constant \(\beta\), feedback can only
\emph{slide} that pole left along the real axis. It cannot split the pole into a
complex pair or cause overshoot.

A proportional--integral (PI) controller is different: its integral term adds a
dynamic state to the feedback path. Oscillation can then arise from the added
dynamics, a higher-order plant, a measurement path, or parasitic poles. Those
effects lie outside the constant-\(\beta\), single-pole model used here.
\end{physicalbox}

\section{Proportional Feedback on a Second-Order System}
Now the plant has two poles. Using the standard form of \cref{ch:linsys} with DC
gain \(A_0\),
\[
A(s)=\frac{A_0\,\omega_0^2}{s^2+2\zeta\omega_0 s+\omega_0^2},
\]
the closed-loop denominator \(1+A(s)\beta=0\) becomes
\[
s^2+2\zeta\omega_0 s+\omega_0^2\bigl(1+A_0\beta\bigr)=0 .
\]
The corresponding closed-loop transfer function is
\[
A_{\mathrm{cl}}(s)=
\frac{A_0\omega_0^2}
{s^2+2\zeta\omega_0s+\omega_0^2(1+A_0\beta)},
\]
so its DC gain is again \(A_{\mathrm{cl}}(0)=A_0/(1+A_0\beta)\).
Compare this with the standard form \(s^2+2\zeta_{\mathrm{cl}}\omega_{\mathrm{cl}}s
+\omega_{\mathrm{cl}}^2\) and the identifications fall out:
\[
\omega_{\mathrm{cl}}=\omega_0\sqrt{1+A_0\beta},
\qquad
\zeta_{\mathrm{cl}}=\frac{\zeta}{\sqrt{1+A_0\beta}},
\qquad
Q_{\mathrm{cl}}=Q\sqrt{1+A_0\beta}.
\]
This is a qualitatively different story from the first-order case. Feedback pushes
the natural frequency \emph{up} but the damping ratio \emph{down}: the system gets
faster and less damped at the same time. Raise \(\beta\) far enough and an initially
sluggish, overdamped plant becomes lightly damped and rings. In pole terms the pair
leaves the real axis and climbs a vertical line toward the \(\jj\omega\) axis, as
\cref{fig:fb-second-order} shows.

\begin{figure}[htbp]
\centering
\includegraphics[width=0.78\textwidth]{feedback_second_order}
\caption{Proportional feedback on a second-order system. \emph{Left:} the poles
break away from the real axis and rise vertically as \(\beta\) increases---the
real part is pinned at \(-\zeta\omega_0\) while the imaginary part grows.
\emph{Right:} the step response develops overshoot and ringing as
\(\zeta_{\mathrm{cl}}\) falls.}
\label{fig:fb-second-order}
\end{figure}

Notice that for this ideal, stable two-pole plant under ordinary negative feedback,
the real part stays at \(-\zeta\omega_0\), so the poles approach the axis but never
cross it: the response rings ever harder yet remains stable. Adding a third pole,
feedback-path dynamics, or delay can change that conclusion, which is where the
frequency view earns its place.
\Cref{sec:threepoles} proves that counting claim once the next chapter has the
tool for it---the proof is pure phase accounting, and it also produces the exact
gain at which a three-pole loop breaks into oscillation.

\section{The Frequency View: Loop Gain and Margins}
The state-space and root-locus pictures ask \emph{where are the closed-loop
poles?} The frequency view asks an equivalent but more practical question:
\emph{how close is the loop gain \(L(\jj\omega)\) to the dangerous value
\(-1\)?} Since the closed-loop poles solve \(1+L(s)=0\), a pole sits exactly on the
\(\jj\omega\) axis when
\[
L(\jj\omega)=-1,
\qquad\text{that is}\qquad
|L(\jj\omega)|=1
\;\;\text{and}\;\;
\angle L(\jj\omega)=-180^\circ
\]
at the same frequency. The virtue of this form is that \(|L|\) and \(\angle L\) are
exactly what a Bode plot of the loop gain shows, so stability becomes something you
read off a measured plot rather than compute from a model.

Two numbers quantify the distance from that condition:
\begin{description}
\item[Gain margin] at the frequency where the phase reaches \(-180^\circ\) (the
\emph{phase crossover}), how far \(|L|\) is \emph{below} \(\SI{0}{\decibel}\).
It answers: how much more gain could I add before oscillating?
\item[Phase margin] at the frequency where \(|L|=1\) (\(\SI{0}{\decibel}\), the
\emph{gain crossover}), how far the phase is \emph{above} \(-180^\circ\). It
answers: how much extra phase lag---an additional negative phase shift---could the
loop tolerate?
\end{description}
For the stable open-loop, negative-feedback setting modeled here, positive margins
mean the closed-loop poles are in the left half-plane. Small margins mean poles near
the \(\jj\omega\) axis: peaked frequency response,
overshoot, long ringing. Zero margin puts poles on the axis---sustained
oscillation. \Cref{fig:fb-margins} shows both margins on one loop-gain Bode plot.

\subsection{One Point, One Plane: the Nyquist View}
\label{sec:nyquist}
The Bode plot separates the loop gain \(L(\jj\omega)\) into magnitude and phase.
A \textbf{Nyquist plot} instead traces that same complex loop gain as frequency
sweeps. In \cref{fig:nyquist}, the curve starts near \(L=K\) at low frequency and
approaches the origin at high frequency. The critical point is \(-1\): reaching it
means \(|L|=1\) and \(\angle L=-180^\circ\) at the same frequency.

This one picture shows why gain and phase margins are complementary. The gain
margin is the radial clearance to \(-1\) along the negative real axis; the phase
margin is the remaining angle from the unit-circle crossing to that axis. The
example uses the same three-pole loop as \cref{fig:fb-margins}, so it connects the
two Bode margins to a single geometric clearance.

\begin{figure}[htbp]
\centering
\includegraphics[width=0.86\textwidth]{nyquist_margins}
\caption{Nyquist plot of the three-pole loop \(L(s)=K/(1+s/\omega_0)^3\) at
\(K=3\)---the same loop, at the same gain, whose Bode margins appear in
\cref{fig:fb-margins}. Both margins are clearances from the single point
\(-1\): the \emph{gain} margin is radial (the locus crosses the negative real
axis at \(-K/8=-0.375\), so the gain could rise \(2.67\times\), or
\SI{8.5}{\decibel}, before reaching \(-1\)), and the \emph{phase} margin is
angular (\(42^\circ\) of arc remain from the unit-circle crossing round to
\(-180^\circ\)). Raising \(K\) scales the whole curve outward until it passes
through \(-1\) at \(K=8\).}
\label{fig:nyquist}
\end{figure}

This chapter needs only that geometric interpretation. \Cref{ch:highorder}
accounts for the phase that can bring a higher-order loop to the critical point,
and \cref{ch:active} applies the same condition to active circuits.

\begin{mistakebox}
\textbf{Two unrelated things are named after Nyquist.} The plot above, and the
\emph{Nyquist sampling theorem}---which says a signal must be sampled at more than
twice its highest frequency, \(f_s>2f_{\max}\), or the frequencies above
\(f_s/2\) (the \emph{Nyquist frequency}) fold back and masquerade as lower ones.
That folding is called aliasing, and it is why a digitizing receiver puts an
anti-aliasing low-pass filter ahead of its converter. \Cref{ch:dsp} develops all of
that properly. The two ideas share only a surname: this chapter's Nyquist plot is
about feedback stability and has nothing to do with sampling rates.
\end{mistakebox}

\begin{figure}[htbp]
\centering
\includegraphics[width=0.86\textwidth]{feedback_margins}
\caption{Gain and phase margin on a loop-gain Bode plot for a three-pole loop.
In the upper magnitude panel, the gain crossover (\(|L|=\SI{0}{\decibel}\)) fixes
the phase margin. In the lower phase panel, the phase crossover
(\(\angle L=-180^\circ\)) fixes the gain margin. Raising the loop gain slides the
magnitude curve up and shrinks both margins toward zero---the onset of
oscillation.}
\label{fig:fb-margins}
\end{figure}

\begin{controlsbox}
\textbf{Two views, one system.} Root locus and Bode margins are the same
statement seen from different sides. The locus shows \emph{where} the poles are as
a parameter varies, which is the natural view when you are choosing that
parameter. The margins show \emph{how much room} the current design has, which is
the natural view when you have measured a loop and want to know if it is safe.
A design usually uses both: the locus to place, the margins to verify.
\end{controlsbox}

\section{Why This Matters in Radio}
Nothing in this chapter required a particular circuit. The results apply to the
linear time-invariant models considered here, with the stated feedback-path and
stability assumptions. Part~III now uses them in circuits:
\begin{itemize}
\item The \textbf{large-loop-gain limit} \(A_{\mathrm{cl}}\to1/\beta\) is why an
op-amp's gain is set by a pair of resistors instead of by the amplifier
(\cref{ch:active}). An ideal op-amp is just a forward path with \(A\to\infty\).
\item The \textbf{gain--bandwidth trade} is the op-amp gain--bandwidth product: the
same invariant \(A_0\omega_a\), renamed.
\item \textbf{Zero margin}---poles parked on the \(\jj\omega\) axis---is not a
failure but a design goal when you want an \emph{oscillator}. The Barkhausen
condition is exactly \(L=-1\).
\item \textbf{Small margins} explain why an active filter that behaves at audio
frequencies misbehaves at RF, and why a voltage regulator can ring.
\end{itemize}
\Cref{ch:active} uses these feedback ideas for op-amps, active filters,
oscillators, and regulators.

\begin{exambox}
The exam does not ask you to compute a margin. It asks which oscillator is most
stable, why an amplifier breaks into oscillation, what negative feedback does to
gain and bandwidth, and what sets an op-amp's gain. All four are the results
above: feedback trades gain for bandwidth, large loop gain makes the closed-loop
behavior depend only on the feedback network, and oscillation is loop gain of
\(-1\).
\end{exambox}
